% =============================================================================
% SYSTEM INSTRUCTION TO LLM:
% This file is a PROMPT TEMPLATE.
% You MUST generate a FULL, FINAL, COMPILABLE LaTeX presentation
% based on this template.
%
% CRITICAL RULES:
% 1. DO NOT change anything in the Preamble.
% 2. You MUST delete all [STYLE REF] frames and the AI INSTRUCTION frame.
% 3. You MUST replace all placeholder texts with real content.
% 4. You MUST preserve the Execushares theme and visual style.
% 5. Keep each slide's body text to roughly 8--10 lines, More than 10 pages is unacceptable.
% 6. NEVER use \includegraphics. No external image files exist. Draw every chart,
%    diagram and table with tikz / pgfplots / booktabs instead.
% 7. Only these packages are available (already loaded in the preamble / theme):
%    Execushares (loads tikz, hyphenat, fontspec + Source Pro fonts),
%    tikz (positioning), pgfplots, amsmath, amssymb, pifont, booktabs, colortbl,
%    multirow. Do NOT add \usepackage lines.
%    Runtime fonts required: Source Serif Pro, Source Sans Pro, Source Code Pro.
% =============================================================================

\documentclass[aspectratio=169,10pt]{beamer}

% -----------------------------------------------------------------------------
% 1. Preamble & Theme (DO NOT CHANGE)
% -----------------------------------------------------------------------------
\usetheme{Execushares}

% Theme already loads tikz; extend libraries and other helpers used in demos
\usetikzlibrary{positioning}
\usepackage{pgfplots}
\pgfplotsset{compat=1.17}
\usepackage{amsmath,amssymb}
\usepackage{pifont}
\usepackage{booktabs}
\usepackage{colortbl}
\usepackage{multirow}

% -----------------------------------------------------------------------------
% 2. Metadata (LLM MUST FILL)
% -----------------------------------------------------------------------------
\title[Short Title]{Full Presentation Title}
% Caller fills these three fields. Leave empty when unknown; never invent placeholders.
\subtitle{}
\author{}
\institute{}
\date{\today}

\begin{document}

% =============================================================================
% FRAME 1: TITLE PAGE
% =============================================================================
\maketitle

% =============================================================================
% FRAME 2: MANUAL TABLE OF CONTENTS
% =============================================================================
\begin{frame}{Agenda}
\large
\renewcommand{\baselinestretch}{1.8}\selectfont

01 \quad First Section Title

02 \quad Second Section Title

03 \quad Third Section Title

04 \quad Fourth Section Title

05 \quad Summary \& Q\&A
\end{frame}

% =============================================================================
% SECTION 1 (EXAMPLE)
% IMPORTANT: define \section{...} before any content frames of that section.
% Execushares also auto-inserts a section splash page via \AtBeginSection.
% =============================================================================
\section{First Section Title}

% =============================================================================
% FRAME 3: SECTION TRANSITION (EXAMPLE: SECTION 1)
% LLM MUST generate one such page before EACH section
% =============================================================================
\begin{frame}{Agenda}
\large
\renewcommand{\baselinestretch}{1.8}\selectfont

\textbf{\textcolor{ExecusharesRed}{01 \quad First Section Title}}

02 \quad Second Section Title

03 \quad Third Section Title

04 \quad Fourth Section Title

05 \quad Summary \& Q\&A
\end{frame}

% =============================================================================
% CONTENT FRAME TEMPLATE (LLM MUST USE THIS STRUCTURE)
% =============================================================================
% Content layout 1: text left, figure right
\begin{frame}{Content Example: Text Left, Figure Right}
\framesubtitle{} % keep empty
\small
\begin{columns}[T]
\column{0.55\textwidth}
Write the core analysis on the left.
Use this area for background, problem definition, reasoning, or methodology.
Keep the text dense and conclusion-oriented; avoid vague filler.

\vspace{0.6em}
\textbf{Key Takeaway:}
State the single main point this slide must convey.

\column{0.40\textwidth}
\centering
% Side panel: FEW categories only (4--6 short labels). Compact width. No xlabel.
% For 12-month / many-label series use the full-width chart on the next example page.
\begin{tikzpicture}
\begin{axis}[
    width=5.4cm,
    height=3.6cm,
    ybar,
    bar width=9pt,
    enlarge x limits=0.15,
    ymin=0, ymax=100,
    symbolic x coords={Q1,Q2,Q3,Q4},
    xtick=data,
    ylabel={Sample index},
    tick label style={font=\tiny},
    label style={font=\scriptsize},
    axis y line=left,
    axis x line=bottom,
    axis line style={ExecusharesRed},
    grid=major,
    grid style={ExecusharesRed!12},
]
\addplot[fill=ExecusharesRed!40, draw=ExecusharesRed] coordinates {
    (Q1,42) (Q2,55) (Q3,61) (Q4,78)
};
\end{axis}
\end{tikzpicture}
\end{columns}
\end{frame}

% Content layout 2: text above, figure below
\begin{frame}{Content Example: Text Above, Figure Below}
\framesubtitle{}
\small

% Upper text region
When a series has many categories (for example 12 months), use this
full-width layout below -- not the narrow side-panel chart.
Keep the prose short so the wide chart remains readable.

\vspace{1em}

% Lower figure: MANY categories (12 months) → MUST be wide + rotated x labels
\centering
\begin{tikzpicture}
\begin{axis}[
    width=11.5cm, height=4.2cm,
    enlarge x limits=0.04,
    x tick label style={rotate=45, anchor=east, font=\tiny},
    ylabel style={font=\scriptsize},
    xlabel style={font=\scriptsize, yshift=-2pt},
    ymin=1900, ymax=2900,
    symbolic x coords={Jan,Feb,Mar,Apr,May,Jun,Jul,Aug,Sep,Oct,Nov,Dec},
    xtick=data,
    ylabel={Sample level},
    xlabel={Month (illustrative)},
    legend style={at={(0.02,0.98)}, anchor=north west, font=\tiny, draw=none, fill=none},
    yticklabel style={font=\tiny},
    axis y line=left,
    axis x line=bottom,
    grid=major,
    grid style={ExecusharesRed!12},
    axis line style={ExecusharesRed},
]
\addplot[
    thick,
    mark=*,
    ExecusharesRed,
    mark options={fill=white, draw=ExecusharesRed, scale=0.7},
] coordinates {
    (Jan,2040) (Feb,2085) (Mar,2160) (Apr,2330) (May,2345) (Jun,2320)
    (Jul,2420) (Aug,2505) (Sep,2630) (Oct,2740) (Nov,2655) (Dec,2620)
};
\addplot[
    thick,
    dashed,
    ExecusharesRed!55,
    mark=square*,
    mark options={fill=white, draw=ExecusharesRed!55, scale=0.55},
] coordinates {
    (Jan,2010) (Feb,2050) (Mar,2120) (Apr,2280) (May,2300) (Jun,2290)
    (Jul,2380) (Aug,2460) (Sep,2580) (Oct,2700) (Nov,2620) (Dec,2590)
};
\legend{Series A (sample), Series B (sample)}
\end{axis}
\end{tikzpicture}

\vspace{0.3em}
\begin{center}
    \tiny \emph{Source: data compiled from the source document.}
\end{center}
\end{frame}


% Content layout 3: figure left, text right
\begin{frame}{Content Example: Figure Left, Text Right}
\framesubtitle{}
\small
\begin{columns}[c] % [c] vertical center; use [T] for top align if text is long

% Left: real TikZ structure diagram (layered pipeline)
\column{0.45\textwidth}
    \centering
    \begin{tikzpicture}[
        font=\scriptsize,
        box/.style={
            draw=ExecusharesRed, fill=ExecusharesRed!12, rounded corners=2pt,
            minimum height=0.62cm, text width=1.35cm, align=center
        },
        layer title/.style={font=\tiny\bfseries, text=ExecusharesRed}
    ]
        % Data layer
        \fill[ExecusharesRed!6, rounded corners=3pt] (-0.25,1.05) rectangle (4.85,2.45);
        \node[layer title, anchor=north west] at (-0.15,2.40) {Data layer};
        \node[box] (s1) at (0.7,1.55) {Sensors};
        \node[box] (s2) at (2.5,1.55) {Ingest};
        \node[box] (s3) at (4.2,1.55) {Store};
        % Decision layer
        \fill[ExecusharesRed!6, rounded corners=3pt] (-0.25,-0.95) rectangle (4.85,0.70);
        \node[layer title, anchor=north west] at (-0.15,0.65) {Decision layer};
        \node[box] (p1) at (0.7,-0.15) {Feature};
        \node[box] (p2) at (2.5,-0.15) {Model};
        \node[box] (p3) at (4.2,-0.15) {Act};
        % Edges
        \draw[->, thick, ExecusharesRed] (s1) -- (s2);
        \draw[->, thick, ExecusharesRed] (s2) -- (s3);
        \draw[->, thick, ExecusharesRed] (s2) -- (p1);
        \draw[->, thick, ExecusharesRed] (p1) -- (p2);
        \draw[->, thick, ExecusharesRed] (p2) -- (p3);
        \draw[->, thick, ExecusharesRed, dashed]
            (p3.south) -- ++(0,-0.28) -| (p1.south)
            node[pos=0.25, below, font=\tiny, text=ExecusharesRed] {feedback};
    \end{tikzpicture}

% Right: analysis region (50% width)
\column{0.50\textwidth}
    \textbf{Object / figure analysis:}

    Interpret the visual on the left. Suitable for:
    \begin{itemize}
        \item Physical structure of equipment
        \item Module-level details of a model
        \item A focused vertical data slice
    \end{itemize}

    \vspace{1em}
    \textbf{Key findings:}
    \begin{enumerate}
        \item \textbf{Observation 1}: The left figure shows...
        \item \textbf{Observation 2}: Relative to expectation, the result...
    \end{enumerate}
\end{columns}
\end{frame}

% =============================================================================
% [STYLE REF] TABLE PAGE (DELETE IN FINAL OUTPUT)
% =============================================================================
\begin{frame}{[STYLE REF] Table Style Example}
\framesubtitle{}
\small
\centering
\renewcommand{\arraystretch}{1.15}
\footnotesize
\rowcolors{2}{ExecusharesRed!10}{white}
\begin{tabular}{p{0.2\textwidth} c c c}
\toprule
Metric & A & B & Verdict \\
\midrule
Item one & 12 & 15 & Room to improve \\
Item two & 20 & 18 & Stable \\
\bottomrule
\end{tabular}
\end{frame}

% =============================================================================
% [STYLE REF] TIKZ FLOW PAGE (DELETE IN FINAL OUTPUT)
% =============================================================================
\begin{frame}{[STYLE REF] Flowchart Style Example}
\framesubtitle{}
\centering
\begin{tikzpicture}[
node distance=0.9cm and 0.7cm,
process/.style={rectangle,draw=ExecusharesRed,fill=ExecusharesRed!10,
minimum height=0.95cm,text width=1.9cm,align=center}
]
\node[process] (a) {Collect};
\node[process,right=of a] (b) {Analyze};
\node[process,right=of b] (c) {Decide};
\node[process,below=of b] (d) {Revise};
\draw[->,thick,ExecusharesRed] (a)--(b);
\draw[->,thick,ExecusharesRed] (b)--(c);
\draw[->,thick,ExecusharesRed] (b)--(d)
    node[midway,right,font=\tiny,text=ExecusharesRed] {gap};
\draw[->,thick,ExecusharesRed,dashed] (d.west) -| (a.south)
    node[near end,left,font=\tiny,text=ExecusharesRed] {retry};
\end{tikzpicture}
\end{frame}

% =============================================================================
% [STYLE REF] TWO-COLUMN ANALYSIS (DELETE IN FINAL OUTPUT)
% =============================================================================
\begin{frame}{[STYLE REF] Two-Column Analysis Example}
\framesubtitle{}
\small
\begin{columns}
\column{0.52\textwidth}
Left side for argument, definition, or derivation.

\column{0.43\textwidth}
Right side for actions, diagrams, or comparisons.
\end{columns}
\end{frame}

% =============================================================================
% AI INSTRUCTION PAGE (DELETE IN FINAL OUTPUT)
% =============================================================================
\begin{frame}[allowframebreaks]{AI GENERATION INSTRUCTIONS}
\framesubtitle{}
You MUST:
\begin{itemize}
\item Delete all [STYLE REF] frames and this instruction frame
\item Replace all placeholders with real content
\item Insert a highlighted TOC transition page before each section
\item Define \texttt{\textbackslash section\{...\}} before the first frame of that section
\item Use columns layout for all content pages
\item Keep information density high (no empty slides)
\item Output FULL LaTeX code only
\end{itemize}
\end{frame}

% =============================================================================
% FINAL PAGE
% =============================================================================
\begin{frame}{Thank You}
\centering
\vspace{1.5em}
{\Large Thank you for your attention.}

\vspace{1.2em}
{\large Questions \& Discussion}

\vspace{1.5em}
{\small \insertauthor \quad\textbar\quad \insertinstitute}
\end{frame}

\end{document}
